\section{Worked Experiments}
\label{sec:experiments}

To ground the pedagogical claims of Section~\ref{sec:pedagogy} in something a reader can
reproduce rather than take on faith, this section specifies six complete circuits, each
built from the mod's own components and each demonstrating a distinct piece of standard
analog-circuits curriculum: a resistive divider, a first-order RC low-pass frequency response,
a second-order RLC resonance, a diode rectifier, a memristor's characteristic pinched
hysteresis loop, and an operational amplifier's open-loop gain rolloff. Two of these
(Experiments~\ref{sec:exp-rc} and~\ref{sec:exp-rlc}) use the AC (small-signal) solver and its
dedicated Bode-plot probe (Section~\ref{sec:ac-analysis}) rather than the time-domain
oscilloscope, since a filter's frequency response is the quantity a Bode plot is specifically
designed to show; the diode rectifier and memristor hysteresis loop remain time-domain
experiments instead, since both are genuinely nonlinear, large-signal phenomena that AC
small-signal analysis cannot represent at all (Section~\ref{sec:ac-analysis} discusses this
limitation directly for both the diode and, more permanently, the memristor). Every circuit
was built directly into the same persistent world used throughout this paper's own
development, laid out as six side-by-side benches on a flattened platform beginning a few
blocks east of the world's spawn point, so a reader with access to that world can walk up to
any of them immediately rather than reconstruct them from a diagram. Because every predicted
value below follows from closed-form analysis of the exact component values specified
(Sections~\ref{sec:circuit-solver} and~\ref{sec:ac-analysis} give the governing equations for
each element in the transient and AC cases respectively), a reader who instead prefers to
build these circuits independently, anywhere, need only reproduce the topology and preset
values given here.

\subsection{A methodological note on presets}

None of the mod's block entities persist their right-click-selected preset (resistance,
capacitance, waveform kind, and so on) across a world save and reload -- confirmed directly
against the source, no block entity in the mod overrides Minecraft's block-entity NBT
read/write methods. Consequently every component below, freshly placed, loads at its
hardcoded default preset regardless of which value the experiment actually calls for, and
reaching the intended value is itself the first step of performing the experiment, exactly
as a student wiring up a real bench supply must first turn its dial to the desired voltage
rather than finding it pre-set. Each experiment below states the default and the exact
number of right-clicks needed to reach the intended value, using the fixed cycle order each
preset wraps through (Section~\ref{sec:circuit-solver} and the appendix give the full preset
tables); the AC Source's amplitude and frequency range, having no preset cycle of their own
at all (Section~\ref{sec:ac-analysis}), are instead given as the exact values to type into
the shift-right-click value editor. A power supply or function generator additionally
remains electrically disconnected -- an open circuit, not a 0~V source -- until it receives
a redstone signal (Section~\ref{sec:circuit-solver}); each bench using one includes a lever
placed directly on top of its source block for exactly this purpose, left off by default so
that flipping it marks a clean $t=0$ for whichever transient the experiment observes. The AC
Source needs no such lever: it is only ever driven during an AC sweep, triggered directly by
the AC probe rather than by continuous per-tick simulation (Section~\ref{sec:architecture}).

\subsection{Experiment 1: basic voltage divider}
\label{sec:exp-divider}

The simplest circuit in this set: a DC source and two resistors in series, tapped at their
midpoint. Built as a single row -- power supply, resistor $R_1$, a wire tap, resistor $R_2$
-- closed into a loop by a plain wire return path and a ground reference, ten blocks east of
spawn. Both resistors are left at their default 100~$\Omega$ preset and the supply at its
default 5~V preset, so this experiment requires no right-clicks at all beyond flipping the
lever.

\textbf{Procedure.} Flip the lever on top of the power supply. Pin the oscilloscope probe to
the wire block between the two resistors (the divider's midpoint) to read its absolute
voltage relative to the ground block elsewhere in the same network.

\textbf{Expected result.} By the standard voltage-divider relation,
\begin{equation}
V_{\text{tap}} = V_{\text{supply}} \cdot \frac{R_2}{R_1 + R_2} = 5~\text{V} \cdot \frac{100}{100+100} = 2.5~\text{V},
\end{equation}
a steady 2.5~V once the lever is flipped, with no transient of any kind (every element here
is purely resistive). A useful follow-up, left to the reader: right-click $R_2$ twice
(100~$\Omega \to 1000~\Omega \to 10\,000~\Omega$) and confirm the tap voltage rises to
$5 \cdot 10\,000/10\,100 \approx 4.95$~V, reproducing the divider relation at a different
ratio without changing the topology at all. A purely resistive network's response is, in
fact, exactly this same flat ratio at every frequency -- pinning the AC probe here instead
(after replacing the power supply with an AC Source, as in Experiments~\ref{sec:exp-rc}
and~\ref{sec:exp-rlc} below) would show a perfectly flat 0~dB, $0^\circ$ Bode plot across the
entire swept range, the trivial special case that motivates why an AC analysis is only
interesting once a reactive element is introduced.

\subsection{Experiment 2: RC low-pass Bode plot}
\label{sec:exp-rc}

An AC Source, a resistor, and a capacitor in series, with the capacitor's far lead tied to
ground -- the same series topology this experiment used in earlier versions of this paper for
a transient step response, but driven and probed here entirely differently, as a frequency
sweep rather than a single time-domain charging curve. Built 26 blocks east of spawn.
Right-click the resistor twice ($100 \to 1000 \to 10\,000\,\Omega$); right-click the
capacitor once ($10 \to 100\,\mu\text{F}$). This gives the same
$\tau = RC = 10\,000 \cdot 100\times10^{-6} = 1$~s time constant as before, corresponding to a
cutoff frequency $f_c = 1/(2\pi\tau) \approx 0.159$~Hz. Shift-right-click the AC Source and
set its frequency range to a minimum of $0.1$~Hz and a maximum of $100$~Hz -- comfortably
bracketing $f_c$ on both sides in log-frequency, so the flat passband, the cutoff itself, and
the asymptotic rolloff are all visible in a single sweep; its amplitude can be left at its
$1$~V default, since a Bode plot's magnitude and phase are both independent of the driving
amplitude entirely.

\textbf{Procedure.} Right-click the AC probe against the AC Source to pin it (a hint
confirms it is pinned), then right-click the capacitor to run the sweep and display the Bode
plot.

\textbf{Expected result.} The standard first-order low-pass transfer function,
\begin{equation}
H(j\omega) = \frac{1}{1 + j\omega RC},
\end{equation}
predicts a response close to flat (about $-1.4$~dB) at the sweep's $0.1$~Hz lower bound, an
exact $-3$~dB with a $-45^\circ$ phase shift at $f_c \approx 0.159$~Hz, and a steep rolloff to
roughly $-56$~dB by the sweep's $100$~Hz upper bound -- a drop of very close to $20$~dB for
every subsequent decade of frequency once well past cutoff, the same asymptotic rate a
first-order filter always exhibits regardless of its specific $R$ and $C$ values. Because
this same $R$ and $C$ combination was previously characterized by its transient step response
(a rise to within $0.7\%$ of $12$~V by $t \approx 5$~s, in an earlier version of this
experiment), a reader who builds both versions can directly compare the two equivalent
descriptions of the identical physical filter: a time constant in the time domain and a
cutoff frequency in the frequency domain, related by nothing more than
$f_c = 1/(2\pi\tau)$.

\subsection{Experiment 3: RLC resonance Bode plot}
\label{sec:exp-rlc}

The same idea as Experiment~\ref{sec:exp-rc}, but with an inductor inserted between the
resistor and capacitor, turning a first-order rolloff into a second-order response capable of
resonant peaking. Built 42 blocks east of spawn. Right-click the resistor \emph{three} times,
cycling $100 \to 1000 \to 10\,000 \to 10~\Omega$ -- the fourth and smallest preset, reached
by wrapping back around past the top of the cycle rather than by cycling downward directly.
Right-click the inductor twice ($0.1 \to 1 \to 5$~H) and the capacitor twice
($10 \to 100 \to 1000\,\mu\text{F}$) -- both to their largest available preset, exactly the
same $R=10\,\Omega$, $L=5$~H, $C=1000\,\mu\text{F}$ combination this experiment used for its
transient ringing response in earlier versions of this paper. Shift-right-click the AC
Source and set its frequency range to a minimum of $0.2$~Hz and a maximum of $20$~Hz, a range
centered (in log-frequency) almost exactly on this circuit's own resonance.

\textbf{Procedure.} Right-click the AC probe against the AC Source, then right-click the
capacitor to run the sweep.

\textbf{Expected result.} With the same undamped natural frequency and damping ratio as
before,
\begin{equation}
\omega_0 = \frac{1}{\sqrt{LC}} \approx 14.14~\text{rad/s}, \qquad
\zeta = \frac{R}{2}\sqrt{\frac{C}{L}} \approx 0.0707,
\end{equation}
the standard second-order low-pass magnitude response peaks at
\begin{equation}
f_{\text{peak}} = \frac{\omega_0\sqrt{1-2\zeta^2}}{2\pi} \approx 2.24~\text{Hz}, \qquad
|H|_{\text{peak}} = \frac{1}{2\zeta\sqrt{1-\zeta^2}} \approx 7.09~\;(\approx 17.0~\text{dB}),
\end{equation}
a pronounced resonant bump rather than a monotonic rolloff, since $\zeta \ll 1/\sqrt{2}$
(the threshold below which any peaking occurs at all). Notably, $f_{\text{peak}}$ here
coincides almost exactly with the damped ringing frequency $f_d \approx 2.24$~Hz this same
$R$, $L$, $C$ combination was shown to ring at in its transient step response in earlier
versions of this experiment -- the same underlying physical resonance manifesting as a peak
in the frequency domain and as visible ringing in the time domain, two views of one
phenomenon rather than two different ones. Past the peak, the response falls off at the
second-order asymptotic rate of $-40$~dB/decade, roughly $-38$~dB below the peak by the
sweep's $20$~Hz upper bound.

\subsection{Experiment 4: half-wave rectifier}
\label{sec:exp-rectifier}

A function generator driving a diode and a load resistor in series, with the resistor's far
lead grounded so it, too, reads an absolute voltage directly. Built 58 blocks east of spawn.
The diode is placed facing \emph{toward} the generator (its anode is the lead facing the
direction it was placed, per Section~\ref{sec:circuit-solver}), so that current flows readily
generator-to-load on the forward half of each cycle and is blocked on the reverse half. No
right-clicks are required at all: the function generator's own default -- a 5~V, 1~Hz sine,
once its lever is flipped -- and the diode's default silicon preset (approximately 0.7~V
forward drop) and the resistor's default 100~$\Omega$ load are exactly the values used. This
experiment deliberately remains a time-domain one rather than an AC sweep: rectification is a
genuinely large-signal, nonlinear phenomenon (the diode's behavior differs qualitatively
between the forward and reverse half-cycles), whereas AC small-signal analysis only ever
linearizes a diode about one fixed operating point (Section~\ref{sec:ac-analysis}) and could
show only a single flat small-signal resistance, not a rectified waveform.

\textbf{Procedure.} Flip the lever. For the clearest comparison, pin the function generator
itself (giving the undistorted input, since a component's trace is the voltage across its
own two leads) as one channel of the time-domain probe and the load resistor (giving the
rectified output) as a second channel, so both waveforms are visible on the same
oscilloscope at once.

\textbf{Expected result.} On each positive half-cycle, once the instantaneous input exceeds
the diode's forward voltage, the output approximately follows the input minus that drop,
peaking around $5 - 0.7 = 4.3$~V; on each negative half-cycle the diode is reverse-biased and
the output reads approximately 0~V. At 1~Hz, ten such rectified pulses appear across the 10~s
history window -- a textbook half-wave-rectified sine, clearly distinguishable from the
smooth input sine on the companion channel.

\subsection{Experiment 5: memristor pinched hysteresis loop}
\label{sec:exp-memristor}

A circuit using the X-Y oscilloscope probe (Section~\ref{sec:architecture}) rather than the
time-domain one: a function generator driving a memristor in series with an ammeter, tracing
the memristor's current against its own voltage. A pinched hysteresis loop through the
origin -- present at any drive frequency, collapsing to a single-valued curve only in the DC
limit -- is the experimental signature that originally motivated the memristor's postulation
and its subsequent physical confirmation (Section~\ref{sec:related-work}). Built 74 blocks
east of spawn, with a frequency module (default 0.5~Hz, requiring no clicks) placed against
the function generator's north face to slow its default 1~Hz down to 0.5~Hz; the generator's
default 5~V amplitude and sine waveform are used unchanged. The memristor is left at its
default $R_{on}=100\,\Omega$, $R_{off}=10\,000\,\Omega$, and $q_{max}=10^{-3}$~C preset
values (all three independently adjustable since Section~\ref{sec:limitations}'s value
editor, but unchanged here). Like Experiment~\ref{sec:exp-rectifier}, this remains a
time-domain experiment rather than an AC sweep: the AC solver's memristor case is, by
deliberate simplification, simply a resistor frozen at whatever resistance it currently
holds, with no frequency dependence of its own at all (Section~\ref{sec:ac-analysis}) -- an
AC sweep here could show only a flat response at that one frozen resistance, never the
hysteresis loop that is the entire point of this experiment.

\textbf{Procedure.} Flip the lever. Pin the memristor \emph{first} with the X-Y probe and
the ammeter \emph{second}: the probe's pinning rule (Section~\ref{sec:architecture}) always
assigns the most recently pinned position to the vertical (Y) axis, so pinning in this order
puts the memristor's voltage on X and the ammeter's current on Y.

\textbf{Expected result.} Unlike Experiments~\ref{sec:exp-divider}
through~\ref{sec:exp-rlc}, this one is genuinely nonlinear -- the memristor's own resistance
depends on the full time history of the current through it, which in turn depends on that
same resistance -- so no closed-form prediction is offered here; the qualitative shape is the
point of the experiment. The current-voltage trace should appear as a pinched loop crossing
through the origin on every cycle (the defining signature of a memristor, as opposed to an
ordinary resistor's straight line or an ordinary reactive element's open ellipse), widest
where the driving voltage is largest and progressively narrowing toward the origin. Because
this mod implements the original, unwindowed HP linear-drift model
(Section~\ref{sec:related-work}), a rough charge-balance estimate for these parameters
(amplitude, frequency, and $q_{\max}$) suggests the accumulated charge over each half-cycle
approaches a substantial fraction of $q_{\max}$, so the loop may visibly flatten as it
approaches its boundary within a half-cycle rather than tracing a smooth figure-eight
throughout -- itself a direct, hands-on illustration of the boundary-pinning limitation that
Biolek et al.'s windowed model was later designed to address, discussed already in
Section~\ref{sec:related-work} and Section~\ref{sec:limitations}.

\subsection{Experiment 6: op-amp open-loop Bode plot}
\label{sec:exp-opamp}

The only circuit in this set using the ideal op-amp, and the only one whose predicted result
comes not from the DC/transient nullor model of Section~\ref{sec:circuit-solver} but from the
two-pole AC model of Section~\ref{sec:ac-analysis} -- deliberately run \emph{open-loop} (with
no feedback network at all), since the entire point is to observe the op-amp's own raw
frequency response directly, rather than a closed-loop response shaped by external
components. Built 83 blocks east of spawn, on a short extension of the same flattened
platform the other five benches sit on: a ground block, then an AC Source (facing east) with
its back lead against that ground; two wire blocks continuing east from the source's front
lead, the second of which turns a corner upward and then east again to arrive directly above
a fourth position; the op-amp itself sits there, placed facing \emph{north} rather than east
(so its non-inverting input, always the up face for a horizontally-facing op-amp per
Section~\ref{sec:architecture}, lands exactly where that overhead wire run terminates) --
its output (north face) is left completely unconnected, probed directly rather than routed
anywhere, and its inverting input (south face, opposite its facing) is tied to a second,
separate ground block placed immediately south of it. No right-clicks at all: the AC
Source's default $1$~V amplitude is used unchanged, and its frequency range is set (via
shift-right-click) to a minimum of $1$~Hz and a maximum of $10$~MHz, spanning several decades
either side of both of the op-amp's fixed pole frequencies.

\textbf{Procedure.} Right-click the AC probe against the AC Source, then right-click the
op-amp block itself (probing a three-terminal component directly reads its output node's
voltage, exactly as the ordinary time-domain probe already does for it) to run the sweep.

\textbf{Expected result.} With the two-pole model's fixed values $A_0 = 10^5$,
$f_{p1} = 20$~Hz, $f_{p2} = 3$~MHz,
\begin{equation}
A(j\omega) = \frac{A_0}{\left(1 + j\dfrac{f}{f_{p1}}\right)\left(1 + j\dfrac{f}{f_{p2}}\right)},
\end{equation}
the sweep's $1$~Hz lower bound reads a gain very close to the full open-loop DC value,
$20\log_{10}(A_0) = 100$~dB, with negligible phase shift; by $f_{p1} = 20$~Hz the gain has
dropped to exactly $97$~dB (the defining $-3$~dB-from-DC property of a pole) with a phase
shift approaching $-45^\circ$; the response then rolls off at the standard single-pole rate
of $-20$~dB/decade until it approaches the second, much higher pole at $f_{p2} = 3$~MHz,
where the rolloff steepens toward $-40$~dB/decade and the phase continues shifting toward
$-180^\circ$ in total. This is, deliberately, the one experiment in this set whose predicted
values cannot be cross-checked against the transient solver at all: the DC/transient model
of Section~\ref{sec:circuit-solver} treats this same op-amp as an ideal, infinite-bandwidth
nullor with no rolloff of any kind, precisely the idealization the AC model exists to move
beyond.
